\documentclass[11pt]{article}
\usepackage{amsmath,amssymb,array,booktabs,geometry}
\geometry{margin=1in}
\newcommand{\half}{\tfrac{1}{2}}
\newcommand{\hZ}{\tfrac{1}{2}\mathbb{Z}}

\title{Index Calculation Debug: \texttt{m007} (Trefoil complement)}
\author{Pipeline verification}
\date{}

\begin{document}
\maketitle

\noindent
\textbf{Manifold:} \texttt{m007} in the SnapPy census (figure-eight knot complement / $4_1$).\\
\textbf{Tetrahedra:} $n=3$. \quad \textbf{Cusps:} $r=1$.\\
\textbf{Basis decomposition:} $1$ hard edge, $1$ easy edge ($n - r = 2$ internal edges total).\\
\textbf{Fugacity variables:} $\eta_1$ (one per hard edge).

%===========================================================
\section{SnapPy Gluing Equations}
%===========================================================

The standard gluing equations are encoded in two integer matrices.
Columns are ordered
$(Z_1, Z_1', Z_1'', \; Z_2, Z_2', Z_2'', \; Z_3, Z_3', Z_3'')$.

\bigskip
\noindent\textbf{Internal edge equations} $C$ ($3\times 9$):
\[
C \;=\;
\begin{pmatrix}
 1 & 0 & 0 & 2 & 0 & 0 & 1 & 0 & 0 \\
 0 & 2 & 1 & 0 & 1 & 2 & 0 & 2 & 1 \\
 1 & 0 & 1 & 0 & 1 & 0 & 1 & 0 & 1
\end{pmatrix}
\]

\noindent\textbf{Cusp equations} $C_{\mathrm{cusp}}$ ($2\times 9$, rows $\mu$, $\lambda$):
\[
C_{\mathrm{cusp}} \;=\;
\begin{pmatrix}
 0 &  1 &  0 & -1 &  0 &  0 &  0 &  0 &  1 \\[3pt]
 2 &  0 &  1 &  0 & -1 & -1 & -1 &  0 & -1
\end{pmatrix}
\]
Row 1 is the meridian $\mu$; row 2 is the longitude $\lambda$.

%===========================================================
\section{Geometric Basis (Hard / Easy Edges)}
%===========================================================

We choose a basis of $n - r = 2$ internal edge equations.
Each row vector is a row of $C$ (same $9$-column format).

\bigskip
\noindent\textbf{Easy edge} (all non-zero entries lie in the $Z_i$ columns, i.e.\ all $Z'$-/$Z''$-entries are zero):
\[
v_{\text{easy}} \;=\; (1,\; 0,\; 0,\quad 2,\; 0,\; 0,\quad 1,\; 0,\; 0)
\]

\noindent\textbf{Hard edge} (has non-zero $Z'$-/$Z''$-entries):
\[
v_{\text{hard}} \;=\; (0,\; 2,\; 1,\quad 0,\; 1,\; 2,\quad 0,\; 2,\; 1)
\]

Because there is one hard edge, the refined index depends on \emph{one} fugacity
variable $\eta_1$, while the easy edge contributes only to the $q$-power.

%===========================================================
\section{Neumann--Zagier Matrix $g_{\mathrm{NZ}}$ and Affine Shift}
%===========================================================

After the symmetry-reduction step the $2n \times 2n = 6\times 6$ matrix
$g_{\mathrm{NZ}}$ is assembled in \emph{block form}
\[
g_{\mathrm{NZ}} \;=\;
\left(\begin{array}{c|c}
  A & B
\end{array}\right),
\qquad
\text{columns: }
(Z_1, Z_2, Z_3 \mid Z_1'', Z_2'', Z_3'').
\]
Rows are ordered: meridian $\mu$; hard-edge charge $Q_{\mathrm{h}}$;
easy-edge charge $Q_{\mathrm{e}}$; longitude/2 $\lambda/2$;
internal momentum $\Gamma_1$ (hard); internal momentum $\Gamma_2$ (easy).

\[
g_{\mathrm{NZ}} \;=\;
\left(\begin{array}{rrr|rrr}
-1 & -1 &  0 & -1 & 0 &  1 \\
-2 & -1 & -2 & -1 & 1 & -1 \\
 1 &  2 &  1 &  0 & 0 &  0 \\\hline
 1 & \half & {-\half} & \half & 0 & {-\half} \\
 0 & -1 &  0 &  0 & 0 &  0 \\
 \half & {-\frac{3}{2}} & 0 & \half & 0 & \half
\end{array}\right)
\]

\medskip
\noindent\textbf{Affine shift.}
The logarithmic branch choices yield integer vectors
$\boldsymbol{\nu}_X, \boldsymbol{\nu}_P \in \hZ^n$:
\[
\boldsymbol{\nu}_X \;=\; (1,\; 3,\; -2),
\qquad
\boldsymbol{\nu}_P \;=\; \left(-\tfrac{1}{2},\; 0,\; 0\right).
\]

\noindent\fbox{%
\parbox{0.9\textwidth}{%
\textbf{Key feature:} $(\boldsymbol{\nu}_P)_1 = -\tfrac{1}{2} \notin \mathbb{Z}$.
This half-integer shift is responsible for the phase selection rule that
forces all odd-$m$ contributions to vanish in the unrefined ($e{=}0$) index.
Note: when the local charges $\widetilde{m}_k, \widetilde{e}_k$ are all
integers, the phase exponent $\varphi$ is \emph{automatically} an integer
(a structural consequence of the NZ matrix).  The \emph{only} integrality
constraint on the summation is that $g_{\mathrm{NZ}}^{-1}\kappa\in\mathbb{Z}^{2n}$.}}

%===========================================================
\section{Inverse Matrix $g_{\mathrm{NZ}}^{-1}$ and Local Charges}
%===========================================================

\subsection{The inverse matrix}

\[
g_{\mathrm{NZ}}^{-1} \;=\;
\left(\begin{array}{rrrrrr}
 \half  &  0 & \half  &  1 &  1 & 0 \\
 0      &  0 &  0     &  0 & -1 & 0 \\
{-\half}&  0 & \half  & -1 &  1 & 0 \\
{-1}    &  0 &{-\half}& -1 & -2 & 1 \\
{-\half}&  1 & \frac{3}{2} & -1 & -1 & 2 \\
 \half  &  0 &  0     &  0 & -2 & 1
\end{array}\right)
\]

\subsection{Master charge vector}

The charge vector entering the summation is
\[
\kappa \;=\;
\begin{pmatrix}
 m \\ 0 \\ 0 \\ e \\ e_1 \\ e_2
\end{pmatrix},
\]
where $m \in \mathbb{Z}$ is the meridian charge, $e \in \mathbb{Z}$ is (half) the
longitude charge, the two zeros fix the position charges of the internal edges
(always zero in the summation), and
$e_1, e_2 \in \tfrac{1}{2}\mathbb{Z}$ are the momenta of the hard and easy
internal edges, respectively. The hard-edge momentum $e_1$ contributes the
fugacity $\eta_1^{2e_1}$ to every term.

\subsection{Local tetrahedron charges}

Set $\Delta = g_{\mathrm{NZ}}^{-1}\,\kappa$.
The local charges split as $(\widetilde{m}_1, \widetilde{m}_2, \widetilde{m}_3,\, \widetilde{e}_1, \widetilde{e}_2, \widetilde{e}_3)$:

\begin{align*}
\widetilde{m}_1 &= \tfrac{m}{2} + e + e_1, \\
\widetilde{m}_2 &= -e_1, \\
\widetilde{m}_3 &= -\tfrac{m}{2} - e + e_1,\\[6pt]
\widetilde{e}_1 &= -m - e - 2e_1 + e_2, \\
\widetilde{e}_2 &= -\tfrac{m}{2} - e - e_1 + 2e_2, \\
\widetilde{e}_3 &= \tfrac{m}{2} - 2e_1 + e_2.
\end{align*}

\subsection{Phase exponent}

Each term in the index carries the factor $(-1)^{\varphi}\,q^{\varphi/2}$ from the
affine shift:
\begin{align*}
\varphi &= \sum_{k=1}^{3} \widetilde{m}_k\,(\boldsymbol{\nu}_P)_k
           - \sum_{k=1}^{3} \widetilde{e}_k\,(\boldsymbol{\nu}_X)_k \\
        &= \widetilde{m}_1\cdot\bigl(-\tfrac{1}{2}\bigr)
           - \widetilde{e}_1\cdot 1
           - \widetilde{e}_2\cdot 3
           - \widetilde{e}_3\cdot(-2).
\end{align*}

Substituting the local charges gives the closed form
\[
\varphi
\;=\;
\frac{13m}{4} + \frac{7e}{2} + \frac{e_1}{2} - 5e_2.
\]

The \emph{only} constraint on the summation is that the local charges
$\widetilde{m}_k, \widetilde{e}_k \in \mathbb{Z}$ (i.e.\ $g_{\mathrm{NZ}}^{-1}\kappa\in\mathbb{Z}^{2n}$).
Whenever this holds, $\varphi\in\mathbb{Z}$ automatically, so no separate
$\varphi\in\mathbb{Z}$ filter is needed or applied.

%===========================================================
\section{Refined Index Formula and Truncated Output}
%===========================================================

\subsection{Index formula}

\[
\mathcal{I}^{\mathrm{ref}}(m,e;\eta_1;q)
\;=\;
\sum_{\substack{e_1,\,e_2\,\in\,\hZ \\ g_{\mathrm{NZ}}^{-1}\kappa\,\in\,\mathbb{Z}^{2n}}}
\eta_1^{2e_1}\,(-1)^{\varphi}\,q^{\varphi/2}
\prod_{k=1}^{3}
\frac{(-1)^{\widetilde{m}_k} q^{\widetilde{m}_k/2}}
     {1 - (-1)^{\widetilde{m}_k+\widetilde{e}_k} q^{(\widetilde{m}_k+\widetilde{e}_k)/2}}.
\]

The product of three tetrahedron factors is
$\displaystyle\prod_{k=1}^3 \mathcal{T}(\widetilde{m}_k, \widetilde{e}_k;q)$,
where the single-tetrahedron building block is
\[
\mathcal{T}(m',e';q)
\;=\;
\frac{(-1)^{m'} q^{m'/2}}{1-(-1)^{m'+e'}q^{(m'+e')/2}}.
\]

\subsection{Truncated expansion at $(m,e)=(0,0)$}

Computed to order $q^6$ (half-integer power cut-off $12$):
\[
\mathcal{I}^{\mathrm{ref}}(0,0;\eta_1;q)
\;=\;
1
- q(2 + \eta_1)
+ q^2(\eta_1^{-1} - 2 - \eta_1)
+ q^3(2\eta_1^{-1} + 4 + 2\eta_1 + \eta_1^2)
+ q^4(\eta_1^{-1} + 9 + 7\eta_1 + 2\eta_1^2)
+ q^5(-4\eta_1^{-1} + 13 + 12\eta_1 + 4\eta_1^2)
+ q^6(-11\eta_1^{-1} + 11\eta_1 + 2\eta_1^2 - \eta_1^3)
+ O(q^7).
\]

\noindent\textbf{Remark:} The coefficient of $q^6$ at $\eta_1^0$ is \emph{zero}
(no constant-$\eta_1$ term appears at order $q^6$).

\bigskip
\noindent The full coefficient table at $(m,e)=(0,0)$:

\begin{center}
\begin{tabular}{cc|r}
\toprule
$q$-power & $\eta_1$-power & coefficient \\
\midrule
$q^0$ & $\eta_1^0$  &  $1$ \\
$q^1$ & $\eta_1^0$  & $-2$ \\
$q^1$ & $\eta_1^1$  & $-1$ \\
$q^2$ & $\eta_1^{-1}$ & $1$ \\
$q^2$ & $\eta_1^0$  & $-2$ \\
$q^2$ & $\eta_1^1$  & $-1$ \\
$q^3$ & $\eta_1^{-1}$ & $2$ \\
$q^3$ & $\eta_1^0$  &  $4$ \\
$q^3$ & $\eta_1^1$  &  $2$ \\
$q^3$ & $\eta_1^2$  &  $1$ \\
$q^4$ & $\eta_1^{-1}$ & $1$ \\
$q^4$ & $\eta_1^0$  &  $9$ \\
$q^4$ & $\eta_1^1$  &  $7$ \\
$q^4$ & $\eta_1^2$  &  $2$ \\
$q^5$ & $\eta_1^{-1}$ & $-4$ \\
$q^5$ & $\eta_1^0$  & $13$ \\
$q^5$ & $\eta_1^1$  & $12$ \\
$q^5$ & $\eta_1^2$  &  $4$ \\
$q^6$ & $\eta_1^{-1}$ & $-11$ \\
$q^6$ & $\eta_1^0$  &  $0$  \quad (absent) \\
$q^6$ & $\eta_1^1$  & $11$ \\
$q^6$ & $\eta_1^2$  &  $2$ \\
$q^6$ & $\eta_1^3$  & $-1$ \\
\bottomrule
\end{tabular}
\end{center}

\end{document}
